%%%%%%%%%%%%%%%%%%%%%%%%%%%%%%%%
% Introduction
%%%%%%%%%%%%%%%%%%%%%%%%%%%%%%%%

\section{Introduction}

The Bernoulli set model~\cite{bernoulliSets} provides a probabilistic framework for \emph{random approximate sets}---sets that approximate an objective set with independent, element-wise errors parameterized by false positive and false negative rates.
That companion paper establishes the axioms, derives the binomial distributions of error counts, and develops the asymptotic normal limits and confidence intervals for the error rates.

This paper addresses the \emph{compositional} question: what happens when Bernoulli sets are combined through set-theoretic operations?
Practical systems rarely use a single approximate set in isolation.
Encrypted search indexes, database query planners, and distributed systems routinely compose approximate sets through union, intersection, complement, and difference.
Without a compositional algebra, each new combination requires a bespoke error analysis.

\paragraph{Contribution.}
We show that Bernoulli sets are \emph{closed} under set-theoretic composition: every operation on Bernoulli sets yields a (possibly higher-order) Bernoulli set whose error rates are closed-form functions of the operand rates and set cardinalities.
Specifically, we derive:
\begin{enumerate}[nosep]
    \item closed-form error rates for complement, union, intersection, and set difference, together with the monoidal structure they induce,
    \item exact formulas for the relational predicates (subset and equality probabilities) on Bernoulli sets, and
    \item an interval arithmetic extension that propagates uncertain rates through set operations, yielding guaranteed bounds on output error rates.
\end{enumerate}

\paragraph{Organization.}
\Cref{sec:preliminaries} restates the definitions and results from~\cite{bernoulliSets} needed in this paper.
\Cref{sec:set_theory} derives error rates for all four set-theoretic operations and establishes closure properties.
\Cref{sec:relational} treats relational predicates (subset, equality) on Bernoulli sets.
\Cref{sec:intervals} develops interval arithmetic for uncertain rate parameters.
